\begin{proof}
We construct a candidate profile $(\tbmp^\dagger, \tbmx^\dagger)$ of the claimed simplified form together with a Bayes-consistent belief system, and verify the two requirements of \cref{def:pbe}. We proceed by backward induction on the tree (leaves to root), showing that at \emph{every} information set $(h_i, v_i)$ the prescribed action maximizes player $i$'s expected utility $\bbU_i$~\eqref{eq:cond-utility} while all other players follow the profile---that is, the profile is sequentially rational. Throughout, the payoff-relevant belief is the Bayes-consistent posterior $\calF(\cdot \mid v_i)$ over $\bmv_{\Desc(i)}$ at every information set (\cref{fact:simplify}), so part~(b) of \cref{def:pbe} holds automatically and we only argue part~(a).

\paragraph{Base case (leaf nodes).}
For any leaf $i \in \calL$, player $i$ has no children, so $\Desc(i) = \varnothing$. Player $i$'s utility from \cref{eq:utility-raw} reduces to:
\[
U_i(h; v_i) \;=\; y_i \cdot (v_i - p_i).
\]
Expanding $y_i = y_{\Parent(i)} \cdot x_i$ via \cref{eq:yi}, the expected utility conditioned on history $h_i$ and action $x_i$ is:
\[
\bbE[U_i \mid h_i, v_i, x_i] \;=\; x_i \cdot y_{\Parent(i)} \cdot (v_i - p_i),
\]
where $y_{\Parent(i)} \in \{0,1\}$ and $p_i \ge 0$ are both determined by $h_i$.

\emph{Case $y_{\Parent(i)} = 0$:} The utility is $0$ regardless of $x_i$, so any $x_i$ is optimal.

\emph{Case $y_{\Parent(i)} = 1$:} The utility equals $x_i(v_i - p_i)$. Since $x_i \in \{0,1\}$, the unique optimum is $x_i = \one\{v_i \ge p_i\}$.

In both cases, the optimal $x_i$ depends only on $(v_i, p_i)$, not on the full history $h_i$. Hence $\tx_i(h_i, v_i) = \tx^\dagger_i(v_i, p_i)$ for some function $\tx^\dagger_i$.

\paragraph{Inductive step.}
Fix a non-leaf node $i$. As inductive hypothesis, assume that for all $m \in \Desc(i)$ the strategies have been simplified---$\tp^\dagger_m$ depends only on $v_m$ and $\tx^\dagger_m$ only on $(v_m, p_m)$---and are moreover \emph{measurable and bounded}; equivalently, each threshold $\ts^\dagger_m$ (\cref{cor:simplify-buy}) is measurable with $\sup_v |\ts^\dagger_m(v)| < \infty$. This holds at the leaves, where $\ts^\dagger_m(v) = v$ is bounded on the compact $\calV_m$. We show below that node $i$ inherits the same properties, closing the induction.

By \cref{fact:simplify} and the Markov property (\cref{asp:markov}), the posterior distribution satisfies $\calF(\bmv_{\Desc(i)} \mid v_i, h_i) = \calF(\bmv_{\Desc(i)} \mid v_i)$. Under the inductive hypothesis, the strategies of all descendants are functions of their own valuations and local prices only, so the expected utility from \cref{eq:utility-raw} expands as:
\begin{align}
&\bbE[U_i \mid h_i, v_i, x_i, \{p_j\}_{j \in \Children(i)}] \nonumber\\
&= x_i \cdot y_{\Parent(i)} \cdot (v_i - p_i) + \bbE_{\bmv_{\Desc(i)} \sim \calF(\cdot|v_i)}\bigg[\sum_{m \in \Desc(i)} \pi(m,i) \cdot p_m \cdot y_m\bigg]. \label{eq:prf:simplify:expand}
\end{align}
For any $m \in \Desc(i)$, the trade indicator decomposes as $y_m = y_i \cdot y_{m|i} = x_i \cdot y_{\Parent(i)} \cdot y_{m|i}$, where $y_{m|i}$ (\cref{eq:y-conditional}) depends only on valuations and strategies within $\Desc(i)$. Substituting into~\eqref{eq:prf:simplify:expand} and factoring out $x_i \cdot y_{\Parent(i)}$:
\begin{equation}
\label{eq:prf:simplify:factor}
\bbE[U_i \mid h_i, v_i, x_i, \{p_j\}] = x_i \cdot y_{\Parent(i)} \cdot \bigg[(v_i - p_i) + \sum_{j \in \Children(i)} G^{(j)}_i(v_i, p_j)\bigg],
\end{equation}
where $G^{(j)}_i(v_i, p_j)$ is defined in \cref{eq:gain-child}. Each $G^{(j)}_i(v_i, p_j)$ depends only on $v_i$ and $p_j$, since the strategies and valuations within $\Subtree(j)$ are independent of $h_i$ by the inductive hypothesis.

\emph{Pricing.} The factor $x_i \cdot y_{\Parent(i)} \ge 0$ does not depend on $\{p_j\}$, so maximizing $\bbE[U_i \mid \cdot]$ over $\{p_j\}$ reduces to maximizing each $G^{(j)}_i(v_i, p_j)$ independently over $p_j \ge 0$. As $G^{(j)}_i(v_i,\cdot)$ depends only on $v_i$ and the strategies within $\Subtree(j)$, the optimizer is a function of $v_i$ alone; we now show it is attained and well-behaved, which both defines $\tp^\dagger_j$ and propagates the inductive hypothesis.

\emph{Attainment and regularity.} Write, from \cref{eq:gain-child},
\[
G^{(j)}_i(v_i, p) = \pi(j,i)\,p\,\Pr[\ts^\dagger_j(v_j) \ge p \mid v_i] \;+\; \bbE_{\bmv_{\Subtree(j)} \sim \calF(\cdot \mid v_i)}\!\big[\one\{\ts^\dagger_j(v_j) \ge p\}\, D\big],
\]
where $D \coloneqq \sum_{m \in \Desc(j)} \pi(m,i)\, \tp^\dagger_m(v_{\Parent(m)})\, y_{m\mid j} \ge 0$ is independent of $p$ and, by the inductive hypothesis (bounded prices) with $\pi(\cdot,\cdot) \le 1$, bounded. Since $\calV_j$ is compact and $\ts^\dagger_j$ bounded, $\bar s_j \coloneqq \sup_{v}\ts^\dagger_j(v) < \infty$, and any $p > \bar s_j$ gives $G^{(j)}_i(v_i,p) = 0$; the search thus reduces to the compact interval $[0, \bar s_j]$.

For each fixed $v_j$, the acceptance indicator $p \mapsto \one\{\ts^\dagger_j(v_j) \ge p\} = \one_{[0,\, \ts^\dagger_j(v_j)]}(p)$ is upper semicontinuous in $p$ (its superlevel sets are closed), thanks to the accept-when-indifferent convention $\{\ts^\dagger_j \ge p\}$. These indicators are uniformly bounded by $1$ and $D$ is bounded, so the reverse Fatou lemma yields, for every $p \in [0, \bar s_j]$, $\limsup_{p' \to p} G^{(j)}_i(v_i, p') \le G^{(j)}_i(v_i, p)$; that is, $G^{(j)}_i(v_i, \cdot)$ is upper semicontinuous on $[0,\bar s_j]$ (the continuous prefactor $\pi(j,i)p$ preserves this). An upper semicontinuous function on a compact set attains its maximum, so $\argmax_{p \ge 0} G^{(j)}_i(v_i, \cdot) \neq \varnothing$; we select its smallest element as $\tp^\dagger_j(v_i)$. Because $G^{(j)}_i$ is jointly measurable in $(v_i,p)$ and upper semicontinuous in $p$, the measurable maximum theorem makes both $\tp^\dagger_j(v_i)$ and $\max_p G^{(j)}_i(v_i,\cdot)$ measurable in $v_i$, and both are bounded ($\tp^\dagger_j \in [0,\bar s_j]$ and $0 \le \max_p G^{(j)}_i \le \bar s_j + \bbE[D]$). Hence $\tp_i(h_i, j, v_i) = \tp^\dagger_j(v_i)$ is a measurable, bounded function of $v_i$ alone.

Crucially, this argument needs only the measurability and boundedness of the $\Subtree(j)$ strategies---which the induction already carries---together with the accept-when-indifferent convention; it never requires their continuity or any joint regularity, so no circularity arises in the backward induction.

\emph{Buying.} After substituting the optimal prices $\{p^*_j(v_i)\}_{j \in \Children(i)}$, define:
\[
\Phi(v_i) \;\coloneqq\; (v_i - p_i) + \sum_{j \in \Children(i)} G^{(j)}_i\big(v_i, p^*_j(v_i)\big).
\]
Then $\bbE[U_i \mid \cdot] = x_i \cdot y_{\Parent(i)} \cdot \Phi(v_i)$. When $y_{\Parent(i)} = 0$, the utility is $0$ for all $x_i$. When $y_{\Parent(i)} = 1$, the optimal choice is $x_i = \one\{\Phi(v_i) \ge 0\}$, which depends only on $(v_i, p_i)$. Hence $\tx_i(h_i, v_i) = \tx^\dagger_i(v_i, p_i)$. The induced threshold $\ts^\dagger_i(v_i) = v_i + \sum_{j \in \Children(i)} \max_{p} G^{(j)}_i(v_i, p)$ (\cref{cor:simplify-buy}) is measurable and bounded on the compact $\calV_i$, since each $\max_p G^{(j)}_i$ is measurable and bounded; thus node $i$ satisfies the inductive hypothesis, closing the induction.

\paragraph{From the construction to a perfect Bayesian equilibrium.}
The backward induction defines, for every player $i$, the simplified strategies $\tp^\dagger_j$ ($j \in \Children(i)$) and $\tx^\dagger_i$. By construction, the prescribed actions maximize player $i$'s expected utility $\bbU_i\big((\tbmp^\dagger, \tbmx^\dagger), \calF(\cdot\mid v_i) \mid h_i, v_i\big)$~\eqref{eq:cond-utility} over player $i$'s own actions at \emph{every} information set $(h_i, v_i)$, while all other players follow $(\tbmp^\dagger, \tbmx^\dagger)$; in particular, this optimality holds for every offered price $p_i$ and every value of $y_{\Parent(i)}$, hence at on- and off-path information sets alike (the root and leaf cases follow from the same computation restricted to pricing or buying alone). Equivalently, for every deviation $(\tp_i, \tx_i)$,
\[
\bbU_i\big((\tbmp^\dagger, \tbmx^\dagger), \calF(\cdot\mid v_i) \mid h_i, v_i\big) \;\ge\; \bbU_i\big((\tp_i, \tbmp^\dagger_{-i}, \tx_i, \tbmx^\dagger_{-i}), \calF(\cdot\mid v_i) \mid h_i, v_i\big),
\]
which is exactly the sequential-rationality requirement~(a) of \cref{def:pbe}. Together with the Bayes-consistency of the beliefs $\calF(\cdot \mid v_i)$ (\cref{fact:simplify}), this shows that $\big((\tbmp^\dagger, \tbmx^\dagger), \mu\big)$ is a perfect Bayesian equilibrium of the simplified form~\eqref{eq:lem:simplify-sell}.
\qed
\end{proof}
